\documentclass{ltjsarticle}
\usepackage{graphicx}
\usepackage{color}
\usepackage{url}
\usepackage{listings}
\usepackage{booktabs}

\begin{document}
\pagestyle{empty}

\begin{center}
\vspace*{1cm}
\large
{\LARGE 卒業研究報告書}\\
\vspace*{0.8cm}
題目\\
\vspace*{1cm}
{\Huge {マルチモーダル学習による薬剤鑑別}}\\
\vspace{3mm}

\vspace*{3cm}
指導教員\\
\vspace*{0.3cm}
{\LARGE 森山 真光  准教授}\\
%\vspace*{-0.3cm}
%\underline{\hspace*{5cm}}\\
\vspace*{3cm}
報告者\\
\vspace*{0.3cm}

{23--1--211--0009}\\
\vspace*{0.3cm}
{\Huge 呉竹 櫂人}\\
% \vspace*{-0.3cm}
% \underline{\hspace*{5cm}}\\
\vspace*{0.5cm}
近畿大学情報学部情報学科\\
\vspace*{2cm}
令和7年4月30日提出\\
\end{center}

\newpage
\normalsize

%\tableofcontents
\clearpage
\begin{center}
{\bf \Large 概要}
\end{center}

錠剤の外観画像と薬剤名・効能・形状などのテキスト情報を組み合わせ，オープンソースLLMであるQwen3.5にマルチモーダル学習させることで，薬剤鑑別システムを構築することを目的とする．提案手法として「RAGによる検索拡張型」と「ファインチューニング型」の2種類を実装し，ChatGPT・Geminiによるゼロショット識別を比較対象として有効性を比較・評価する．

\newpage

\tableofcontents
\newpage
\setcounter{page}{1}
\pagestyle{plain}

\section{序論}\label{sec1}
\subsection{本研究の背景}

薬剤誤投与は医療安全上の重大な問題であり，患者への身体的危害や医療コストの増大を招く．
医療現場では，バーコードが付与されていない薬剤や一包化調剤など，外観のみから薬剤を鑑別しなければならない場面が存在する．

薬剤鑑別を自動化する研究はこれまでも行われてきた．
Heoら\cite{heo2023pill}は，YOLOv5による刻印文字検出とResNetによる形状・色認識を組み合わせ，言語モデルで刻印文字を補正するパイプライン型識別システムを提案し，韓国の薬剤データセットでTop-1精度85.6\%を達成した．
しかし，本手法は錠剤表面の刻印を主な手がかりとするため，刻印のないカプセルや素錠には適用できないという根本的な制約がある．

Rádliら\cite{radli2024metric}は，錠剤画像の特徴と薬剤説明テキストから抽出した単語埋め込みを統合するマルチモーダルメトリクス学習フレームワークを提案し，Top-1精度93.56\%を達成した．
しかし，評価対象が112クラスという小規模なデータセットに限定されており，日本語薬剤への応用や大規模化への対応が示されていない．

近年は大規模言語モデル（LLM）を薬剤識別に用いる研究も登場している．
Menahaら\cite{gemini2024pill}は，Google GeminiをLLMとして採用し，807枚の薬剤画像に対してTop-1精度95.04\%を達成している．
しかし，商用クローズドソースモデルを使用しているため，患者の薬剤画像を外部APIに送信するプライバシーリスクがあり，特定薬剤に対するファインチューニングも行えない．

錠剤の外観画像と薬剤名・形状・効能などのテキスト情報を同時に入力できるマルチモーダルLLMを用いれば，刻印がない薬剤であっても色・形・大きさなどの視覚的特徴とテキスト情報を組み合わせて識別できると考えられる．
さらにオープンソースモデルであれば患者データを外部に送信せずローカルで動作させることができ，ファインチューニングによって特定薬剤への適応も可能である．
以上の課題から，（1）刻印に依存しないマルチモーダルな識別，（2）プライバシーを保持したローカル動作，（3）新薬追加への柔軟な対応，の3点を同時に満たす手法が求められている．

\subsection{本研究の目的}

本研究では，オープンソースのマルチモーダルLLMであるQwen3.5を用いて，薬剤鑑別システムを構築する．
提案手法として以下の2種類を実装し，比較評価する．

\begin{enumerate}
  \item \textbf{提案手法①　Qwen3.5 + RAG（検索拡張生成）}：薬剤データベースをあらかじめベクトル化し，入力に類似する薬剤情報を自動検索してモデルに渡すことで識別精度を高める手法．新薬追加時はデータベースの更新のみで再学習が不要である．
  \item \textbf{提案手法②　Qwen3.5 + ファインチューニング（FT）}：実薬データでQwen3.5を追加学習し，薬剤鑑別に特化したモデルを構築する手法．
\end{enumerate}

比較対象として，ChatGPT・GeminiによるゼロショットAPI識別を用い，提案手法の有効性を客観的に示す．

\subsection{本報告書の構成}

第2章では関連研究を整理し，本研究の位置づけを明確にする．
第3章では提案手法の詳細および実験計画を述べる．
第4章では実験結果と考察を示す．
第5章では結論と今後の課題を述べる．

\newpage

\section{関連記述・関連研究}\label{sec2}

本研究では以下の3本の先行研究を参考文献とし，各手法の概要と限界を整理する．

\subsection{パイプライン型マルチモーダル識別手法}

Heoら\cite{heo2023pill}は，薬剤識別を「認識ステップ」と「検索ステップ」の2段階に分けたパイプラインシステムを提案した．
YOLOv5で錠剤表面の刻印文字を検出し，ResNet-32で形状・色・形態（錠剤/カプセル）を認識する．
さらに，RNNベースの文字レベル言語モデルで刻印文字を補正し，薬剤データベースとの照合で最終識別を行う．
韓国の薬剤データセット（MFDS）においてTop-1精度85.6\%，米国データセット（NLM）において74.5\%を達成している．

\textbf{問題点}
\begin{itemize}
  \item 刻印のないカプセルや素錠には適用できない．
  \item 各モジュールが独立しており，誤差が積み上がる．
  \item 日本語薬剤への対応が示されていない．
  \item 新薬追加時に識別モジュールの再学習が必要となる場合がある．
\end{itemize}

\subsection{マルチモーダルメトリクス学習による識別手法}

Rádliら\cite{radli2024metric}は，画像特徴と薬剤説明テキストの単語埋め込みを統合するメトリクス学習フレームワークを提案した．
EfficientNetV2をバックボーンとしてRGB・LBP・輪郭・テクスチャの4ストリームで視覚特徴を抽出し，テキスト情報から動的マージントリプレット損失（DMTL）のマージンを決定する．
112クラス・4,480枚で構成するOGYEIv2データセットにてTop-1精度93.56\%を達成している．

\textbf{問題点}
\begin{itemize}
  \item 評価対象が112クラスの小規模なデータセットに限定されており，大規模化への対応が未検証である．
  \item ハンガリー語データセットであり，日本語薬剤への応用が示されていない．
  \item LLMを使用していないため，端から端まで一体的な学習ができない．
  \item 新薬追加のたびにメトリクス学習の再学習が必要である．
\end{itemize}

\subsection{LLMによるゼロショット識別手法}

Menahaら\cite{gemini2024pill}は，Google Gemini LLMを用いて錠剤画像から薬剤情報を提供するシステムを提案した．
画像をバイナリ値に変換してオントロジーベースのプロンプトエンジニアリングとともにGeminiに入力し，テキストおよび音声形式で薬剤説明を出力する．
807枚の薬剤画像に対してTop-1精度95.04\%を達成している．

\textbf{問題点}
\begin{itemize}
  \item 商用クローズドソースモデルであり，特定薬剤へのファインチューニングが不可能である．
  \item 患者の薬剤画像を外部APIに送信するため，プライバシー上のリスクがある．
  \item 利用コストが高く，病院・薬局での継続的な運用が困難である．
  \item ゼロショットであるため，外観が類似した薬剤の細かい識別に限界がある．
\end{itemize}

\subsection{先行研究の整理と本研究の位置づけ}

表\ref{tab:related}に先行研究の比較を示す．

\begin{table}[h]
\centering
\caption{先行研究の比較}
\label{tab:related}
\begin{tabular}{lccccc}
\toprule
手法 & テキスト利用 & FT可能 & プライバシー & 新薬追加 \\
\midrule
Heo (2023) \cite{heo2023pill}           & 刻印のみ & 可 & 低リスク & 再学習要 \\
Rádli (2024) \cite{radli2024metric}     & テキスト埋め込み & 可 & 低リスク & 再学習要 \\
Gemini LLM \cite{gemini2024pill}        & あり & 不可 & 高リスク & 不要 \\
\textbf{提案手法① RAG（本研究）}        & あり & 不要 & 低リスク & DB更新のみ \\
\textbf{提案手法② FT（本研究）}         & あり & 可 & 低リスク & 再チューニング要 \\
\bottomrule
\end{tabular}
\end{table}

以上の問題点を踏まえ，本研究ではオープンソースのマルチモーダルLLMを用いることでファインチューニングを可能にし，プライバシーリスクのない薬剤鑑別システムを構築する．

\newpage

\section{研究内容}\label{sec3}

\subsection{使用モデル：Qwen3.5}

本研究ではオープンソースのマルチモーダルLLMであるQwen3.5を使用する．
選定理由を以下に示す．

\begin{itemize}
  \item オープンソースであるため無償でファインチューニング可能であり，患者データを外部APIに送信する必要がないためプライバシーリスクがない．
  \item 画像とテキストを同時に入力して一括学習できるため，複雑な処理工程が不要である．
  \item ChatGPT・Geminiとの精度比較により，提案手法の優位性を客観的に示せる．
\end{itemize}

\subsection{提案手法①：Qwen3.5 + RAG}

薬剤データベースをあらかじめベクトル化しておき，入力された画像・テキストに類似する薬剤情報を自動検索してQwen3.5に渡すことで識別精度を高める．

\textbf{強み}
\begin{itemize}
  \item 新薬追加時はデータベースを更新するだけでよく，再学習が不要である．
  \item 少ないデータ量でも運用できる．
  \item 「この薬に類似している」という識別根拠を提示できる．
\end{itemize}

\textbf{課題}
\begin{itemize}
  \item 検索精度がそのまま識別精度のボトルネックになる．
  \item 識別のたびにデータベース検索が走るため，処理に時間がかかる．
\end{itemize}

\subsection{提案手法②：Qwen3.5 + ファインチューニング}

実薬データ（錠剤画像＋薬剤テキスト）でQwen3.5を追加学習し，薬剤鑑別に特化したモデルを構築する．

\textbf{強み}
\begin{itemize}
  \item 薬剤鑑別に特化した学習による高い識別精度が期待できる．
  \item 推論が速く，完全にローカル環境で動作する．
  \item 外部APIを使用しないためプライバシーリスクがない．
\end{itemize}

\textbf{課題}
\begin{itemize}
  \item 新薬追加のたびに再学習が必要となる．
  \item ある程度の量の学習データが必要である．
\end{itemize}

\subsection{実験計画}

表\ref{tab:schedule}に実験スケジュールを示す．

\begin{table}[h]
\centering
\caption{実験スケジュール}
\label{tab:schedule}
\begin{tabular}{cll}
\toprule
フェーズ & 時期 & 内容 \\
\midrule
Phase 1 & 4〜5月  & 論文調査・実験計画の確定．参考文献3本の精読と課題整理．研究テーマの最終決定． \\
Phase 2 & 6〜7月  & 環境構築とベースライン実験．ChatGPT・Geminiのゼロショット評価を実施． \\
Phase 3 & 7〜9月  & 提案手法①RAGシステムの構築と評価．提案手法②ファインチューニングと評価． \\
Phase 4 & 9〜10月 & 結果分析・考察．全手法の比較評価まとめ．Ablation Study・誤識別ケースの原因分析． \\
Phase 5 & 10〜1月 & 論文執筆・発表準備．卒業論文の執筆と最終発表スライド作成． \\
\bottomrule
\end{tabular}
\end{table}

\newpage

\section{結果・考察}\label{sec4}

（実験実施後に記述予定）

\newpage

\section{結論・今後の課題}\label{sec5}

（実験実施後に記述予定）

\newpage

\nocite{*}
\bibliographystyle{plain}
\bibliography{reference}

\end{document}
